\section{방법}

\subsection{설계}\label{sec:design}
경고 신호(있음/없음)를 참가자 내 요인으로 둔 단일 요인 설계다.
한 참가자가 두 조건을 모두 수행하고, 조건 순서는 참가자마다 시행 단위로 무작위 배정했다.
이렇게 하면 기기의 고정 입력 지연과 개인의 기본 반응속도가 두 조건에 똑같이 들어가, 조건 간 차이에서 상쇄된다 \citep{pronk2020}.
표~\ref{tab:design}에 바꾼 조건과 고정한 조건을 정리했다.

\begin{table}[h]
\centering\small
\caption{실험 설계: 바꾼 조건과 고정한 조건}\label{tab:design}
\begin{tabular}{p{0.24\linewidth}p{0.68\linewidth}}
\toprule
\multicolumn{2}{l}{\textbf{바꾼 조건(독립변수)}}\\
신호 있음 & 목표 500\,ms 전, 화면 안쪽 테두리(두께 10\,px, 밝은 회백색)가 100\,ms 점등. 1인당 20회\\
신호 없음 & 같은 시점에 화면 변화 없음. 1인당 20회\\
\midrule
\multicolumn{2}{l}{\textbf{고정한 조건(통제변수)}}\\
경고 시점 $\to$ 목표 & 두 조건 모두 500\,ms (신호 없음 조건에서는 보이지 않는 가상 시점)\\
시행 시작 $\to$ 경고 시점 & 1000--2500\,ms 균등분포에서 시행마다 추출, 두 조건 같은 분포\\
목표 자극 & 화면 중앙 흰 원, 지름 120\,px (CSS 픽셀)\\
응답 & 화면 아무 곳 터치·클릭(\texttt{pointerdown}), 또는 Space/Enter 키(\texttt{keydown})\\
반응 창 & 목표 후 1500\,ms. 넘으면 놓침으로 기록\\
시행 수 & 연습 6회(조건별 3회, 분석 제외) + 본 시행 40회\\
\bottomrule
\end{tabular}
\end{table}

\subsection{참가자}
연구자의 지인과, 연구자가 실험 링크를 올린 온라인 커뮤니티의 이용자를 대상으로 편의 표집했다.
수집 기간은 링크를 처음 공유한 날부터 7일로 미리 정했다.
실제로는 \CollectStart{}부터 \CollectEnd{}까지 모았고, 이 기간에 본 실험을 끝낸 세션은 \Nrecruited{}개였다.
목표 인원은 20명 이상으로 정했다.
이 숫자는 검정력 분석이 아니라 1주 안에 모을 수 있는 범위를 기준으로 정했으며, 그 한계는 \ref{sec:limits}절에 적었다.
이름·연락처 등 개인을 식별하는 정보는 받지 않았다.
연령대는 선택 응답으로 받았고, 참가자는 시작 화면에서 기록되는 항목을 확인한 뒤 동의했다.

\subsection{장치와 자극}
실험 페이지는 React와 Vite로 만들었고 Vercel에 배포했다(\ExperimentURL).
참가자는 각자의 휴대폰, 태블릿, PC 브라우저로 참여했다.
자극 제시와 시간 측정은 React 렌더링을 거치지 않고 DOM을 직접 바꾸는 별도 모듈에서 처리했다.
목표 자극은 정해진 시각에 가장 가까운 \texttt{requestAnimationFrame} 콜백에서 그렸고, 그 시점의 \texttt{performance.now()}를 제시 시각으로 기록했다.
반응 시각은 입력 이벤트의 \texttt{event.timeStamp}로 기록했다.
\texttt{click}은 손을 뗄 때 발생해 불필요한 지연이 섞이므로, 누르는 순간 발생하는 \texttt{pointerdown}을 썼다.
난수는 \texttt{crypto.getRandomValues}로 만들었다.

\subsection{절차}
참가자는 안내와 동의 화면 다음에 방법 설명을 읽고 연습 6회를 수행했다.
매 시행은 화면 가운데 작은 점이 나타나며 시작된다.
1000--2500\,ms 뒤의 경고 시점에, 신호 있음 조건이면 테두리가 100\,ms 동안 밝아졌다.
그로부터 500\,ms 뒤 점이 사라지고 흰 원이 나타났다.
참가자는 원이 보이면 최대한 빨리 화면을 누르거나 키를 눌렀다.
원이 나오기 전에 누르면 ``너무 일러요''라는 안내와 함께 그 시행을 조기 반응으로 기록했다.
반응 뒤에는 그 시행의 반응시간을 600\,ms 동안 보여 주고, 500\,ms 쉰 뒤 다음 시행으로 넘어갔다.
본 시행 40회를 마치면 개인 결과를 보여 주고 기록을 서버로 보냈다.
전체 소요 시간은 약 3분이었다.

\subsection{측정값}
종속변수는 단순 반응시간, 즉 반응 시각에서 목표 제시 시각을 뺀 값(ms)이다.
시행마다 조건, 계획한 대기시간, 실제 경고--목표 간격, 상태(정상·조기 반응·놓침), 반응시간, 입력 방식, 시행 도중 탭 전환 여부를 기록했다.
세션마다 앱 버전, 같은 브라우저의 재참여 여부, 포인터 종류, 화면 크기, 추정 프레임 간격을 기록했다.
참가자 1명의 기록은 JSON 파일 하나로 저장했다.
모든 원자료와 정리된 CSV는 재현 패키지에 들어 있다.

\subsection{측정 도구 점검}
수집 전에 브라우저 안의 ``로봇 참가자''로 페이지를 점검했다.
로봇은 원이 뜨면 정확히 250\,ms 뒤에 누르도록 만들었다.
데스크톱 Chromium에서 46행(연습 6, 본 시행 40)이 설계대로 기록되었다.
실제 경고--목표 간격은 499.4--500.3\,ms(설계 500\,ms)였고, 기록된 반응시간은 250.2--250.5\,ms였다.
일부러 넣은 조기 반응 1건은 반응시간 없이 조기 반응으로 기록되었다.
즉 소프트웨어 측정 오차는 1\,ms 미만이었다.
다만 로봇은 이벤트를 소프트웨어로 발생시키므로, 이 점검에는 터치스크린·키보드 하드웨어의 지연이 포함되지 않는다.

\subsection{사전 등록한 분석}
분석 규칙은 데이터 수집 전에 문서와 코드로 고정했고, 공개 Git 저장소의 커밋 기록으로 시점을 남겼다.
\textbf{시행 제외}: 연습 시행, 조기 반응, 놓침, 시행 도중 탭·앱 전환이 감지된 시행, 반응시간 100\,ms 미만.
\textbf{참가자 제외}(순서대로 적용):
같은 브라우저의 두 번째 이후 참여,
본 시행을 모두 마치지 않은 세션,
유효 시행 30개 미만이거나 어느 한 조건의 유효 시행이 10개 미만인 경우.
\textbf{분석 단위}: 참가자별·조건별 유효 반응시간의 중앙값과, 그 차이 $d = $ 중앙값(신호 있음) $-$ 중앙값(신호 없음).
반응시간 분포는 오른쪽 꼬리가 길어, 평균 대신 이상값의 영향을 덜 받는 중앙값을 썼다 \citep{ratcliff1993}.
\textbf{확증 검정}: $d$에 대한 단측 Wilcoxon 부호순위 검정(대립가설: 중앙값 $<0$), $\alpha = .05$.
$p<.05$이고 $d$의 중앙값이 음수이면 ``가설이 지지되었다'', 그 외에는 ``지지되지 않았다''로 판정한다.
\textbf{보조 분석}: 단측 대응 $t$검정, 효과크기 $d_z$, 평균 차이의 95\% 신뢰구간,
$d$ 중앙값의 부트스트랩 95\% 신뢰구간(재표집 10{,}000회, 시드 20261030), 부호 검정, 민감도 분석.
\textbf{탐색 분석}(확증 아님): 입력 방식별 $d$, 조건별 조기 반응 건수.
분석은 Python(pandas, SciPy)으로 했고, 결과 수치와 그림은 분석 스크립트가 원자료에서 바로 만들어 이 문서에 넣었다.
